\documentclass[tikz,border=8pt]{standalone}
\usepackage{tikz}
\usepackage{amsmath}
\usepackage{amssymb}
\usepackage{pgfplots}
\pgfplotsset{compat=1.18}
\usepackage{ctex}
\usetikzlibrary{calc,positioning,arrows.meta,matrix}

% LC 20 括号匹配（Valid Parentheses）
% 字符串 "({[]})" 处理过程，栈竖立，5 个快照

\begin{document}
\begin{tikzpicture}[
  every node/.style={font=\small},
  stackcell/.style={draw=gray!70, rectangle,
    minimum width=0.85cm, minimum height=0.65cm, thick, fill=blue!8},
  pushcell/.style={draw=green!60!black, rectangle,
    minimum width=0.85cm, minimum height=0.65cm, very thick, fill=green!10},
  popcell/.style={draw=red!60!black, rectangle,
    minimum width=0.85cm, minimum height=0.65cm, very thick, fill=red!10},
  matchlabel/.style={font=\scriptsize,green!60!black},
  curchar/.style={draw=orange!70,rectangle,rounded corners=2pt,
    minimum width=0.7cm,minimum height=0.6cm,thick,fill=orange!10},
  ptr/.style={->,>=Stealth,thick},
]

%% ── 标题 ──────────────────────────────────────────────────
\node[font=\large\bfseries] at (7.5, 1.0)
  {括号匹配（LC 20 Valid Parentheses）};
\node[font=\small,gray] at (7.5, 0.4)
  {字符串 \texttt{"(\{[]\})"} ：遇左括号入栈，遇右括号与栈顶匹配则出栈};

%% ── 字符串展示（顶部） ──────────────────────────────────
\def\chars{{"(", "\{", "[", "]", "\}", ")"}}

\foreach \i in {0,1,2,3,4,5} {
  \pgfmathsetmacro\c{\chars[\i]}
  \node[draw=gray!60,rectangle,minimum width=0.7cm,minimum height=0.6cm,
        fill=white,thick,font=\ttfamily\small]
    at (4.0+\i*0.85, -0.2) {\c};
  \node[font=\tiny,gray] at (4.0+\i*0.85, -0.65) {\i};
}

%% ── 辅助：画竖直栈（从底往上），共 cols，栈底 y=ya ──
% 参数：列 x，底 y，内容列表（从底到顶），高亮最顶

%% ─────────────────────────────────────────────────────────
%% 快照列布局：每列 x 位置
%% snap 0,1,2,3,4 共 5 个，每列宽约 2.4cm
%% ─────────────────────────────────────────────────────────

\def\snapxA{1.2}
\def\snapxB{3.8}
\def\snapxC{6.4}
\def\snapxD{9.0}
\def\snapxE{11.6}
\def\stacktop{-1.6} % y coordinate of stack top label position
\def\stackbase{-6.0} % y of stack bottom

%% ── 快照 0：处理 '('，入栈 ──
\node[font=\scriptsize\bfseries,blue!70] at (\snapxA, \stacktop+0.35)
  {处理 \texttt{(}};
\node[font=\tiny,gray] at (\snapxA, \stacktop+0.0) {idx=0，入栈};

% 当前字符高亮
\node[curchar,font=\ttfamily] at (4.0+0*0.85, -0.2) {(};

% 栈：只有 ( 在位置 0
\node[pushcell,font=\ttfamily] (a0) at (\snapxA, -2.5) {(};
% 栈底线
\draw[thick,gray!60] (\snapxA-0.55, -3.0) -- (\snapxA+0.55, -3.0);
\node[font=\tiny,gray] at (\snapxA, -3.25) {栈底};
\node[font=\tiny,green!60!black,anchor=west] at (\snapxA+0.5, -2.5) {$\leftarrow$ top};

%% ── 快照 1：处理 '\{'，入栈 ──
\node[font=\scriptsize\bfseries,blue!70] at (\snapxB, \stacktop+0.35)
  {处理 \texttt{\{} };
\node[font=\tiny,gray] at (\snapxB, \stacktop+0.0) {idx=1，入栈};

% 栈从底到顶：( 在底，{ 在顶（高亮）
\node[stackcell,font=\ttfamily] (b0) at (\snapxB, -3.1) {(};
\node[pushcell,font=\ttfamily]  (b1) at (\snapxB, -2.4) {\{};
\draw[thick,gray!60] (\snapxB-0.55, -3.6) -- (\snapxB+0.55, -3.6);
\node[font=\tiny,gray] at (\snapxB, -3.85) {栈底};
\node[font=\tiny,green!60!black,anchor=west] at (\snapxB+0.5, -2.4) {$\leftarrow$ top};

%% ── 快照 2：处理 '['，入栈 ──
\node[font=\scriptsize\bfseries,blue!70] at (\snapxC, \stacktop+0.35)
  {处理 \texttt{[}};
\node[font=\tiny,gray] at (\snapxC, \stacktop+0.0) {idx=2，入栈};

\node[stackcell,font=\ttfamily] (c0) at (\snapxC, -3.8) {(};
\node[stackcell,font=\ttfamily] (c1) at (\snapxC, -3.1) {\{};
\node[pushcell,font=\ttfamily]  (c2) at (\snapxC, -2.4) {[};
\draw[thick,gray!60] (\snapxC-0.55, -4.3) -- (\snapxC+0.55, -4.3);
\node[font=\tiny,gray] at (\snapxC, -4.55) {栈底};
\node[font=\tiny,green!60!black,anchor=west] at (\snapxC+0.5, -2.4) {$\leftarrow$ top};

%% ── 快照 3：处理 ']'，与栈顶 '[' 匹配 → 出栈 ──
\node[font=\scriptsize\bfseries,blue!70] at (\snapxD, \stacktop+0.35)
  {处理 \texttt{]}};
\node[font=\tiny,gray] at (\snapxD, \stacktop+0.0) {idx=3，匹配出栈};

\node[stackcell,font=\ttfamily] (d0) at (\snapxD, -3.1) {(};
\node[stackcell,font=\ttfamily] (d1) at (\snapxD, -2.4) {\{};
% 原栈顶 [ 已匹配弹出（用虚线框表示已移除）
\node[popcell,font=\ttfamily,dashed] (d2) at (\snapxD, -1.7) {[};
\draw[very thick,red!60] (\snapxD-0.3,-1.55) -- (\snapxD+0.3,-1.85);
\draw[very thick,red!60] (\snapxD+0.3,-1.55) -- (\snapxD-0.3,-1.85);

\draw[thick,gray!60] (\snapxD-0.55, -3.6) -- (\snapxD+0.55, -3.6);
\node[font=\tiny,gray] at (\snapxD, -3.85) {栈底};
\node[font=\tiny,green!60!black,anchor=west] at (\snapxD+0.5, -2.4) {$\leftarrow$ top};
\node[font=\tiny,red!70!black,anchor=west] at (\snapxD+0.5, -1.7) {弹出};
\node[matchlabel,anchor=south] at (\snapxD, -1.2) {\texttt{[} 与 \texttt{]} 匹配};

%% ── 快照 4：处理 '\}'，匹配出栈；处理 ')'，匹配出栈 → 栈空=true ──
\node[font=\scriptsize\bfseries,blue!70] at (\snapxE, \stacktop+0.35)
  {处理 \texttt{\}} \texttt{)}};
\node[font=\tiny,gray] at (\snapxE, \stacktop+0.0) {idx=4,5 两次出栈};

% 两个都弹出了，栈空
\node[popcell,font=\ttfamily,dashed] (e1) at (\snapxE, -3.1) {\{};
\draw[very thick,red!60] (\snapxE-0.3,-2.95) -- (\snapxE+0.3,-3.25);
\draw[very thick,red!60] (\snapxE+0.3,-2.95) -- (\snapxE-0.3,-3.25);

\node[popcell,font=\ttfamily,dashed] (e0) at (\snapxE, -3.8) {(};
\draw[very thick,red!60] (\snapxE-0.3,-3.65) -- (\snapxE+0.3,-3.95);
\draw[very thick,red!60] (\snapxE+0.3,-3.65) -- (\snapxE-0.3,-3.95);

\draw[thick,gray!60] (\snapxE-0.55, -4.3) -- (\snapxE+0.55, -4.3);
\node[font=\tiny,gray] at (\snapxE, -4.55) {栈底};

% 结果
\node[draw=green!60!black,fill=green!8,rounded corners=3pt,
      font=\small\bfseries,inner sep=5pt]
  at (\snapxE, -1.7) {\textcolor{green!50!black}{栈空 $\Rightarrow$ true}};

%% ── 分隔线 ──────────────────────────────────────────────
\draw[gray!25,dashed] (0.0, -6.5) -- (14.0, -6.5);

%% ── 算法说明 ─────────────────────────────────────────────
\node[draw=gray!50,fill=white,rounded corners=3pt,font=\small,
      inner sep=8pt,align=left]
  at (7.5, -7.6)
  {\textbf{括号匹配算法（O(n) 时间，O(n) 空间）：}\\[4pt]
   遇左括号 \texttt{( [ \{} $\to$ 入栈\\[2pt]
   遇右括号 \texttt{) ] \}} $\to$ 检查栈顶是否为对应左括号：\\
   \quad 匹配 $\to$ 出栈；\quad 不匹配或栈空 $\to$ 返回 false\\[2pt]
   全部处理后：栈空 $\to$ 返回 true；否则返回 false};

\end{tikzpicture}
\end{document}
